\section{The Geometry of Decay: The Residual Power Law}

The most striking dynamic validation of the $E_8$-Persistence Theory is the strict correlation between a particle's \textbf{Geometric Residual} ($R$) and its \textbf{Mean Lifetime} ($\tau$).

In the $E_8$ framework, a particle is defined as a resonant knot in the lattice. If the knot is perfect ($R=0$), the vacuum cannot "grip" the structure to dismantle it. However, if the geometric formula yields a mass slightly mismatched from the lattice ideal ($M_{exp} \neq M_{geo}$), this \textbf{Residual} ($R$) represents a ``Geometric Frustration,'' a loose thread in the topology.

We posit the \textbf{Residual Hypothesis}: \textit{The magnitude of the geometric error ($R$) defines the effective Phase Space volume available for decay.}

\subsection{The Empirical Regression}

To test this hypothesis, we analyzed the fundamental unstable particles identified in the spectral scan. We excluded composite hadrons (whose lifetimes are complicated by binding energies) and stable particles, isolating the ``Fundamental Instability'' of the lattice.

\textbf{Dataset:} The blind scan identified the following residuals ($|M_{exp} - M_{calc}|$) and lifetimes:

\begin{table}[h]
\centering
\begin{tabular}{lrr}
\toprule
\textbf{Particle} & \textbf{Residual ($R$)} & \textbf{Lifetime ($\tau$)} \\ \midrule
\textbf{Muon}     & $0.09 \, m_e$           & $2.2 \times 10^{-6}$ s \\
\textbf{Strange}  & $0.06 \, m_e$           & $1.2 \times 10^{-8}$ s \\
\textbf{Charm}    & $0.50 \, m_e$           & $1.1 \times 10^{-12}$ s \\
\textbf{Tau}      & $2.168 \, m_e$          & $2.9 \times 10^{-13}$ s \\
\textbf{Bottom}   & $3.48 \, m_e$           & $1.5 \times 10^{-12}$ s \\
\textbf{W Boson}  & $61.98 \, m_e$          & $3.0 \times 10^{-25}$ s \\
\textbf{Top}      & $137.54 \, m_e$         & $5.0 \times 10^{-25}$ s \\ \bottomrule
\end{tabular}
\end{table}

\textbf{Regression Results:} Plotting $\log_{10}(\tau)$ against $\log_{10}(R)$ yields a strict Power Law relation: 
\[ \tau \propto |R|^{\beta} \]

\begin{itemize}
    \item \textbf{Calculated Slope ($\beta$):} \textbf{$-5.49$}
    \item \textbf{Correlation ($R^2$):} \textbf{$0.91$}
\end{itemize}

\subsection{Physical Interpretation: Recovering the Sargent Rule}

The calculated slope of $\beta \approx -5.49$ is physically profound. In standard Quantum Field Theory, the decay rate ($\Gamma$) of a particle via the Weak Interaction is governed by \textbf{Sargent's Law}, which scales with the fifth power of the mass: 
\[ \Gamma \propto G_F^2 m^5 \implies \tau \propto m^{-5} \]

The $E_8$ spectral scan reveals that the \textbf{Geometric Residual ($R$)} acts as the effective kinematic mass driving the decay.

\begin{enumerate}
    \item \textbf{The Geometry of Time:} The universe does not decay based on total mass; it decays based on \textbf{Geometric Error}. The ``Mass'' term in the Sargent Rule is replaced by the ``Residual'' term in the Persistence framework.
    \item \textbf{Phase Space Dominance:} The slope of $-5.49$ (close to the theoretical $-5.0$) confirms that the decay mechanism is dominated by the available phase space volume defined by the mismatch. The slight steepening ($-5.49$ vs $-5.0$) reflects the \textbf{Entropic Multiplier} ($N_{terms}$), where particles with larger residuals also tend to possess higher geometric complexity.
\end{enumerate}

\subsection{Stability and the ``LOCKED'' State}

This power law explains the bifurcation of the universe into ``Eternal'' and ``Ephemeral'' matter:

\begin{itemize}
    \item \textbf{The Unstable ($\tau \to 0$):} Particles like the Top Quark have massive geometric residuals ($R \approx 137 \, m_e$). The vacuum impedance creates a massive drag force ($\tau \propto 137^{-5.49}$), resulting in instantaneous yoctosecond decay.
    \item \textbf{The Stable ($\tau \to \infty$):} The Electron and Proton possess geometric residuals below the \textbf{Persistence Margin} ($PM \approx 10^{-6} \, m_e$).
    \begin{itemize}
        \item Plugging $R < 10^{-6}$ into the power law yields a lifetime exceeding the age of the universe.
        \item They are \textbf{Geometrically Locked}. Because the vacuum has a finite resolution, it cannot resolve the error required to trigger the decay process.
    \end{itemize}
\end{itemize}

\textbf{Conclusion:} Time is not a fundamental variable for particles. It is a function of \textbf{Geometric Fit}. A particle with a perfect integer fit to the $E_8$ lattice experiences zero time; it is eternal. A particle with a poor fit experiences rapid time; it is erased by the vacuum's error-correction protocols.